\section{The breaking point: a reverse stress test on flood severity}\label{sec:stress}

Question~3 turns the chain into a tool. Supervisory climate exercises run forward: they fix a scenario and compute the loss \citep{ecb2022climate,fed2024pilot,boe2022results}. A reverse stress test runs backward: it fixes the loss that matters, a breach of the capital threshold, and solves for the scenario \citep{grundke2018macroeconomic,glasserman2015stress}. For floods, \citet{blickle2025bad} compute the disaster severity that would fail a single-county US bank, and \citet{caloia2021flood} describe their forward exercise on Dutch banks' geocoded collateral, run through increasingly extreme scenarios, as in essence a reverse stress test; the published version \citep{caloia2026floods} uses fixed scenarios and reports non-linear capital depletion. This section formalises the inversion of equation~\eqref{eq:sstar} for a scaling of the flood footprint, with valuation behaviour and relief as explicit inputs, and shows which inputs it is sensitive to.

\subsection{The inversion}

The inversion is tested end to end on the synthetic panel of Appendix~\ref{app:code}, where the truth is known: regions are provinces, flood exposure $x_r$ comes from the 2011 flood shares of the collateral sheets, the flood sensitivity of default $\beta_{\mathrm{pd}}$ is the annualised window effect from the event study, and the loss-given-default parameters take their planted values, with illustrative capital (a starting common equity ratio of \rqThreeCetBefore\%, a baseline default probability of \rqThreePdZero\% and a threshold of \rqThreeKappa\%). With the planted $\beta_s = \rqThreeBetaTrue$ the true breaking point is $S^* = \rqThreeSTrue$ and the estimated $\hat\beta_s = \rqThreeBetaHat$ (\rqThreeBetaSe) recovers \rqThreeSHat; propagating the estimation uncertainty and a uniform prior on $\lambda_B$ over $[0.1, 0.5]$ gives a median of \rqThreeSMedian{} with a 5--95\% range of [\rqThreeSLow, \rqThreeSHigh] over \rqThreeDraws{} draws. The inversion is stable, monotone (checked on a grid before bisection) and recovers the truth to within the sampling error of its inputs. This is a check that the code finds the root of the function that generated the data; it says nothing about any real bank.

\subsection{Does it matter whether valuations price hazard?}

\begin{table}[t]
\centering\footnotesize
\caption{Loss given default and the breaking point under valuations that ignore or price hazard}\label{tab:lambda}
% Generated by 04_analysis/src/paper_numbers.py. Do not edit.
\begin{adjustbox}{max width=\textwidth}
\begin{tabular}{llrrrrrr}
\toprule
$\delta$ & LTV & \multicolumn{2}{c}{LGD, flooded region (\%)} & Change (pp) & \multicolumn{2}{c}{$S^*$} & Change (\%) \\
\cmidrule(lr){3-4}\cmidrule(lr){6-7}
 & & $\lambda_B = 0$ & $\lambda_B = 1$ & & $\lambda_B = 0$ & $\lambda_B = 1$ & \\
\midrule
\multirow{4}{*}{0.023} & 0.7 & 8.2 & 6.0 & 2.2 & 3.26 & 6.23 & 91 \\
 & 0.8 & 19.6 & 17.8 & 1.9 & 2.42 & 2.52 & 4 \\
 & 0.9 & 28.6 & 26.9 & 1.7 & 1.90 & 1.97 & 4 \\
 & 1.0 & 35.7 & 34.2 & 1.5 & 1.57 & 1.62 & 3 \\
\midrule
\multirow{4}{*}{0.046} & 0.7 & 10.3 & 6.0 & 4.3 & 3.11 & 6.23 & 100 \\
 & 0.8 & 21.5 & 17.8 & 3.8 & 2.31 & 2.52 & 9 \\
 & 0.9 & 30.3 & 26.9 & 3.4 & 1.83 & 1.97 & 8 \\
 & 1.0 & 37.2 & 34.2 & 3.0 & 1.52 & 1.62 & 7 \\
\midrule
\multirow{4}{*}{0.092} & 0.7 & 14.6 & 6.0 & 8.6 & 2.83 & 6.23 & 120 \\
 & 0.8 & 25.3 & 17.8 & 7.6 & 2.12 & 2.52 & 19 \\
 & 0.9 & 33.6 & 26.9 & 6.7 & 1.69 & 1.97 & 16 \\
 & 1.0 & 40.3 & 34.2 & 6.1 & 1.42 & 1.62 & 14 \\
\bottomrule
\end{tabular}
\end{adjustbox}

\tabnote{Illustrative two-region book: \lamWFlooded\% of exposure in a flood-prone region with hazard $h = 1$ and 2011-type flood intensity $x = \lamXFlooded$, \lamWOther\% elsewhere with $h = 0$ and $x = \lamXOther$; baseline default probability \lamPdZero\%, flood sensitivity of default $\beta_{\mathrm{pd}} = \lamBetaPd$, damage share $\phi = \lamPhiBase$, recovery haircut $c = \lamHaircutBase$, starting common equity ratio \lamCetBefore\% and threshold \lamKappa\%. $\delta$ is the full-information flood discount and LTV the book loan-to-value ratio at origination, applied to the appraisal whatever $\lambda_B$. $\lambda_B = 0$: collateral values ignore hazard, so the book LTV understates true leverage; $\lambda_B = 1$: values price it fully. LGD from equation~\eqref{eq:lgd} at $S = 1$; $S^*$ from equation~\eqref{eq:sstar}. Source: \texttt{rq3\_lambda\_sensitivity.csv}.}
\end{table}

The parameter that Question~1 bears on, and that the link from $\lambda_T$ to $\lambda_B$ carries into the stress test, enters loss given default through the kink in equation~\eqref{eq:lgd}. It does so under an assumption that must be stated: loans are written at the same book loan-to-value ratio against the appraisal whatever $\lambda_B$, so a bank whose appraisals ignore hazard has lent more against the same parcel. The $\lambda_B$ effect on $S^*$ is therefore a true-leverage effect created at origination, not an effect of valuation on realised losses. If loan amounts were held fixed instead, exposure at default would not depend on $\lambda_B$, $S^*$ on realised losses would be identical under $\lambda_B = 0$ and $1$, and $\lambda_B$ would move only the reported loan-to-value ratio and the timing of staging and expected credit loss.

Under the origination assumption, Table~\ref{tab:lambda} shows how much $\lambda_B$ matters. At a book loan-to-value ratio of 0.8 and a full-information discount of 4.6\%, a bank whose valuations ignore hazard has a loss given default on flooded collateral \lamLgdEightChange{} percentage points higher than one whose valuations price it, and its breaking point is \lamSStarEightZero{} rather than \lamSStarEightOne, a \lamSStarEightChange\% difference; at a ratio of 1.0 the difference is \lamSStarOneChange\%. Halving the discount to 2.3\% gives differences of \lamLoEightChange\% and \lamLoOneChange\% at ratios of 0.8 and 1.0; doubling it to 9.2\% gives \lamHiEightChange\% and \lamHiOneChange\%. At a ratio of 0.7 the breaking point moves from \lamSStarLowLtvZero{} to \lamSStarLowLtvOne, by \lamSStarLowLtvChange\%, and the figure is the same \lamLowLtvOneAll{} for every $\delta$: it is a kink effect. At that leverage the priced valuation keeps the flooded loans below the kink until a flood about six times the 2011 footprint, so the headline depends on where the kink sits, which is set by the assumed haircut, damage share and intensity rather than by the discount. Appendix Table~\ref{tab:lambdagrid} varies the haircut over 0.2--0.4 and the damage share over 0.1--0.3 at $\delta = 4.6\%$: at a ratio of 0.8 the $\lambda_B$ effect ranges from \lamGridEightMin\% to \lamGridEightMax\%, and at 0.7 from \lamGridLowMin\% to \lamGridLowMax\% in the cells where the book breaches the threshold at all, with no breach within ten times the 2011 footprint in \lamGridLowNoBreach{} of the nine cells. The three-figure results at low leverage are therefore artefacts of the kink's position; the 7--9\% results at ordinary leverage are robust to the haircut and damage share.

Two conclusions follow for the design of the project, both conditional on this illustrative book. Valuation pricing moves the breaking point by 7--9\% at ordinary loan-to-value ratios, and by much more at low ratios, where the result depends on the assumed damage share and haircut; that is enough to make $\lambda_B$ worth measuring, which is why loan-level appraisals are part of the data request, and not enough to call Question~1 a precondition for the credit questions. And the sensitivity is greatest for the loans that look safest on the books, those at low loan-to-value ratios, which is where a valuation that ignores hazard does the most to hide risk. Only the $\lambda_B$ half of prediction P5 is illustrated here; the concentration half, moving a share of exposure into the most flood-exposed region, is implemented in the code but not examined, because the portfolio weights available (loans by booking province, with Bangkok at \bkkShare\%) place corporate loans at head offices rather than at their collateral and would make any concentration figure a statement about booking practice.

\subsection{What the empirical version needs}

The empirical reverse stress test for a Thai bank, or for the system, needs four inputs, of which the paper so far supplies the framework for all and the data for one. Portfolio weights by province come from the public provincial loan series and are stable (Section~\ref{sec:aggregates}), with the booking-province caveat just stated. Flood exposure $x_r$ by province comes from the 2011 flood shares of Section~\ref{sec:cross}, with the urban under-detection caveat, or from the World Bank's district-level mapping \citep{worldbank2026ccdrfin}. The flood sensitivity of default and the collateral-damage share need the loan-level estimates of Question~2, or published analogues such as the default response of \citet{albertazzi2026physical} and the relief-censored response of \citet{chitchumnong2025physical} as bounds. And the capitalisation parameter needs Question~1 for the official valuer and loan-level appraisals for $\lambda_B$. With these, the exercise would report $S^*$ in units of the 2011 footprint, with and without relief (prediction P6), under the 7\% and 8\% thresholds. The capital notification provides for a countercyclical buffer of 0--2.5\% at the Bank of Thailand's discretion and did not require one when issued \citep{bot2012capital}; I found no public announcement activating it as of 4 November 2026, so the baseline sets it to zero and 2.5\% is a sensitivity. The 1-in-50-year classification of 2011 in the World Bank note gives $S^*$ a return-period interpretation.
